\chapter{Musculoskeletal X-Ray}

\section*{What Is This Test?}
Musculoskeletal X-rays image bones and joints to detect fractures, dislocations, arthritis, infections, and tumors. They're the first-line imaging for bone and joint problems.

\section*{Alternative Names}
MSK X-ray, bone X-ray, extremity X-ray, skeletal radiograph, orthopedic X-ray


\section*{Principle}
Musculoskeletal X-ray uses radiographic imaging with at least two orthogonal views (AP and lateral) of the bone or joint. The X-ray beam is attenuated differently by cortical bone, trabecular bone, soft tissue, and fat. Fractures appear as radiolucent lines interrupting the cortical bone continuity.

\section*{History}
X-ray imaging of bones and joints was the first clinical application of radiology, beginning in 1896. It remains the primary imaging modality for fracture detection and joint assessment.

\section*{Why This Test Is Ordered}
Your doctor orders an X-ray of a bone or joint when you have pain, swelling, or an injury to that area. It answers the question: do you have a fracture, arthritis, dislocation, or bone tumor?

\section*{Is This a Primary or Secondary Test}
A primary test. X-ray is the first imaging test for bone and joint problems.

\section*{How This Test Is Performed}
The affected body part is positioned against the X-ray detector. Usually two views are taken at different angles (front and side) to provide a complete picture. The test takes about ten to fifteen minutes.

\section*{What Preparation Is Needed}
Remove any jewelry or metal objects near the area being X-rayed. No other preparation is needed.

\section*{Contraindications and Precautions}
Pregnancy is not an absolute contraindication; the examination should be justified and optimized. Inform the team if you are pregnant or may be pregnant. There are no other specific contraindications for a musculoskeletal X-ray.

\section*{Risks}
A musculoskeletal X-ray uses a small amount of ionizing radiation. The dose is very low and localized to the area being imaged. Inform the team if you are pregnant. There are no other significant risks.

\section*{What Happens After the Test}
Results are typically available within thirty minutes. You may be placed in a splint or cast while waiting for results if a fracture is suspected.

\section*{Understanding Your Results}

\section*{Normal Results}
A normal result shows normal bone alignment, no fractures, no joint abnormalities, and normal soft tissues.
\clearpage
